\subsection{Stronger global normalization changes the interpretation}
\label{sec:globalresults}
A bounded follow-up applies signed Clifford extraction and fusion across
commuting intervening rotations to all 12 existing settings, without
changing constructors or readouts. Its input is the compiled IR; its
remaining Clifford frame is retained unless exactly identity up to
phase. In the two-layer ChebyshevRx examples, unitary-preserving
normalization alone reduces trainable Pauli-angle tasks from 4 to 3
and from 8 to 7. These reductions occur in \emph{both} the boundary and
transport arms (Table~\ref{tab:global}). They must not be credited to
the kernel certificate. A multi-qubit Pauli rotation is one angle task,
not necessarily one native physical gate.

\begin{table}[htbp]
\caption{Two-layer ChebyshevRx after global normalization.
Coefficient rank is the exact linear rank of the retained angle
coefficients in the displayed source-feature basis, not a proof of
minimal kernel dimension. B and T denote the earlier boundary and
transport arms.}
\label{tab:global}
\begin{tabular}{lrrrr}
& \multicolumn{2}{c}{2 qubits} & \multicolumn{2}{c}{4 qubits}\\
& B & T & B & T\\
\toprule
Symbols & 7 & 5 & 15 & 13\\
Angle tasks (old) & 4 & 4 & 8 & 8\\
Angle tasks (global) & 3 & 3 & 7 & 7\\
Coefficient rank & 6 & 4 & 14 & 12\\
\bottomrule
\end{tabular}
\end{table}

Across these 12 settings, no additional angle-task reduction appears
when comparing the globally normalized boundary and transport arms.
All source-to-normal-form feature, unitary, kernel and derivative
probes pass; Appendix~\ref{app:global} specifies the independent
rewrite checker and numerical scope. This is a new local structural
experiment, not a rerun of the Linux timing study. In particular,
the historical timing percentages below are \emph{not} speedups over
this stronger global-normalization baseline. The new representation
can absorb the representative internal explanation into standard
normalization followed by the readout-isometry test; exclusivity over
such a composition is not established.

The zero angle-task increment has a simple sufficient explanation.
For a fixed fused representation with angles
$\alpha_r(x,\theta)=u_r(x,\theta)+b_r(\theta)$, deleting certified
input-independent terms $b_r$ cannot remove a trainable-angle task
if $u_r$ still has a nonzero retained-parameter dependence. If this
holds for every trainable rotation, and gate grouping and measurement
bases are unchanged, the per-angle two-shift schedule keeps exactly
the same task count. This is an accounting statement for that schedule,
not a lower bound on all possible kernel evaluators. It applies to the
surviving product-angle terms in the two representative circuits.
